% GENERATED FILE. Edit canonical lessons, then run npm run generate:books.
% canonical-source-hash: fnv1a64:f9661ec1d9abd774
% canonical-lessons: ZH-C105-maojin, ZH-C105-yashua, ZH-C105-yagao, ZH-C105-feizao, ZH-C105-jingzi

\chapter{In the Bathroom}
\label{ch:zh-in-the-bathroom}
\hlchaptermodality{105}

\begin{chapteropening}
\textbf{By the end of this chapter:} \emph{I can name a towel, a toothbrush, toothpaste, soap and a mirror.}

Complete the last lesson of chapter 105: I can name a towel, a toothbrush, toothpaste, soap and a mirror.
\end{chapteropening}

\section[máojīn]{máojīn — towel}

\label{lesson:ZH-C105-maojin}



\begin{quote}
Before the new one: say the Mandarin for paper, then the Mandarin for dictionary.
\end{quote}

\subsection*{You'll want to know: máojīn}
\emph{máojīn} — \textquotedblleft{}towel\textquotedblright{}.

\begin{quote}
\emph{wǒ de máojīn hěn gānjìng} — My towel is clean.
\end{quote}

Its characters come later; for now it is a word for the ear.

\begin{cousinweb}
\emph{máojīn} is built from máo and jīn, both inherited from Middle Chinese, \emph{maw} and \emph{kin} (Guangyun, in Baxter's transcription); the dictionary glosses the first as hair and fur and explains the second as a cloth worn on the head. How the pair came to mean a towel is not settled.
\end{cousinweb}

\begin{grammarlens}[title={What you've built}]
One more word for everyday things.
\end{grammarlens}

\subsection*{Your turn}
\begin{itemize}
\raggedright
  \item \emph{Say it:} \emph{máojīn}
  \item \emph{Say it:} \emph{máojīn}, once more
  \item \emph{Say it:} \emph{wǒ de máojīn hěn gānjìng} — My towel is clean.
  \item \emph{Recall:} say the Mandarin for paper, then the Mandarin for dictionary, then say \emph{máojīn} again
\end{itemize}

\begin{tcolorbox}[breakable,colback=teal!4,colframe=teal!35!black,title={Before you move on}]
What does máojīn mean? (\textquotedblleft{}Towel\textquotedblright{}.) Say it once more.
\end{tcolorbox}
\section[yáshuā]{yáshuā — toothbrush}

\label{lesson:ZH-C105-yashua}



\begin{quote}
Before the new one: say the Mandarin for dictionary, then the Mandarin for towel.
\end{quote}

\subsection*{You'll want to know: yáshuā}
\emph{yáshuā} — \textquotedblleft{}toothbrush\textquotedblright{}.

\begin{quote}
\emph{wǒ de yáshuā zài nǎr} — Where's my toothbrush?
\end{quote}

Its characters come later; for now it is a word for the ear.

\begin{cousinweb}
\emph{yáshuā} is built from yá 'tooth' and shuā 'brush', both inherited from Middle Chinese, \emph{ngae} and \emph{srwaet} (Guangyun, in Baxter's transcription): the two parts of shuāyá, the other way round, make the thing.
\end{cousinweb}

\begin{grammarlens}[title={What you've built}]
One more word for everyday things.
\end{grammarlens}

\subsection*{Your turn}
\begin{itemize}
\raggedright
  \item \emph{Say it:} \emph{yáshuā}
  \item \emph{Say it:} \emph{yáshuā}, once more
  \item \emph{Say it:} \emph{wǒ de yáshuā zài nǎr} — Where's my toothbrush?
  \item \emph{Recall:} say the Mandarin for dictionary, then the Mandarin for towel, then say \emph{yáshuā} again
\end{itemize}

\begin{tcolorbox}[breakable,colback=teal!4,colframe=teal!35!black,title={Before you move on}]
What does yáshuā mean? (\textquotedblleft{}Toothbrush\textquotedblright{}.) Say it once more.
\end{tcolorbox}
\section[yágāo]{yágāo — toothpaste}

\label{lesson:ZH-C105-yagao}



\begin{quote}
Before the new one: say the Mandarin for towel, then the Mandarin for toothbrush.
\end{quote}

\subsection*{You'll want to know: yágāo}
\emph{yágāo} — \textquotedblleft{}toothpaste\textquotedblright{}.

\begin{quote}
\emph{méiyǒu yágāo le} — We're out of toothpaste.
\end{quote}

Its characters come later; for now it is a word for the ear.

\begin{cousinweb}
\emph{yágāo} is built from yá 'tooth' and gāo 'paste', both inherited from Middle Chinese, \emph{ngae} and \emph{kaw} (Guangyun, in Baxter's transcription); the dictionary glosses the second as fat or grease.
\end{cousinweb}

\begin{grammarlens}[title={What you've built}]
One more word for everyday things.
\end{grammarlens}

\subsection*{Your turn}
\begin{itemize}
\raggedright
  \item \emph{Say it:} \emph{yágāo}
  \item \emph{Say it:} \emph{yágāo}, once more
  \item \emph{Say it:} \emph{méiyǒu yágāo le} — We're out of toothpaste.
  \item \emph{Recall:} say the Mandarin for towel, then the Mandarin for toothbrush, then say \emph{yágāo} again
\end{itemize}

\begin{tcolorbox}[breakable,colback=teal!4,colframe=teal!35!black,title={Before you move on}]
What does yágāo mean? (\textquotedblleft{}Toothpaste\textquotedblright{}.) Say it once more.
\end{tcolorbox}
\section[féizào]{féizào — soap}

\label{lesson:ZH-C105-feizao}



\begin{quote}
Before the new one: say the Mandarin for toothbrush, then the Mandarin for toothpaste.
\end{quote}

\subsection*{You'll want to know: féizào}
\emph{féizào} — \textquotedblleft{}soap\textquotedblright{}.

\begin{quote}
\emph{yòng féizào xǐ shǒu} — Wash your hands with soap.
\end{quote}

Its characters come later; for now it is a word for the ear.

\begin{cousinweb}
\emph{féizào} is built from féi 'fat' and zào, both inherited from Middle Chinese, \emph{bj+j} and \emph{dzawX} (Guangyun, in Baxter's transcription). The dictionary glosses the first 'fleshy', and the second as black silk, among other senses; how the pair came to mean soap is not settled.
\end{cousinweb}

\begin{grammarlens}[title={What you've built}]
One more word for everyday things.
\end{grammarlens}

\subsection*{Your turn}
\begin{itemize}
\raggedright
  \item \emph{Say it:} \emph{féizào}
  \item \emph{Say it:} \emph{féizào}, once more
  \item \emph{Say it:} \emph{yòng féizào xǐ shǒu} — Wash your hands with soap.
  \item \emph{Recall:} say the Mandarin for toothbrush, then the Mandarin for toothpaste, then say \emph{féizào} again
\end{itemize}

\begin{tcolorbox}[breakable,colback=teal!4,colframe=teal!35!black,title={Before you move on}]
What does féizào mean? (\textquotedblleft{}Soap\textquotedblright{}.) Say it once more.
\end{tcolorbox}
\section[jìngzi]{jìngzi — mirror}

\label{lesson:ZH-C105-jingzi}



\begin{quote}
Before the new one: say the Mandarin for toothpaste, then the Mandarin for soap.
\end{quote}

\subsection*{You'll want to know: jìngzi}
\emph{jìngzi} — \textquotedblleft{}mirror\textquotedblright{}.

\begin{quote}
\emph{xǐshǒujiān lǐ yǒu jìngzi ma} — Is there a mirror in the toilets?
\end{quote}

Its characters come later; for now it is a word for the ear.

\begin{cousinweb}
\emph{jìngzi} is built from jìng 'mirror' and the noun ending zi, both inherited from Middle Chinese, \emph{kjaengH} and \emph{tsiX} (Guangyun, in Baxter's transcription); the dictionary's entry tells of a mirror sent as a gift to a king of old.
\end{cousinweb}

\begin{grammarlens}[title={What you've built}]
One more word for everyday things.
\end{grammarlens}

\subsection*{Your turn}
\begin{itemize}
\raggedright
  \item \emph{Say it:} \emph{jìngzi}
  \item \emph{Say it:} \emph{jìngzi}, once more
  \item \emph{Say it:} \emph{xǐshǒujiān lǐ yǒu jìngzi ma} — Is there a mirror in the toilets?
  \item \emph{Recall:} say the Mandarin for toothpaste, then the Mandarin for soap, then say \emph{jìngzi} again
\end{itemize}

\begin{tcolorbox}[breakable,colback=teal!4,colframe=teal!35!black,title={Before you move on}]
What does jìngzi mean? (\textquotedblleft{}Mirror\textquotedblright{}.) Say it once more.
\end{tcolorbox}
